% Public source snapshot. See sources/README.md for the project's stricter
% independent-provenance and licensing boundary.
\documentclass{article}
\usepackage{graphicx}
\usepackage{amsmath}
\usepackage{amssymb}
\usepackage{listings}
\usepackage[margin=1in]{geometry}
\usepackage{xcolor}
\usepackage{tikz}
\usetikzlibrary{shapes, arrows, positioning, calc}
\usepackage{tcolorbox}

% Define colors for syntax highlighting
\definecolor{codegreen}{rgb}{0,0.6,0}
\definecolor{codegray}{rgb}{0.5,0.5,0.5}
\definecolor{codepurple}{rgb}{0.58,0,0.82}
\definecolor{backcolour}{rgb}{0.95,0.95,0.92}

% Python style for highlighting
\lstdefinestyle{pythonstyle}{
    backgroundcolor=\color{backcolour},
    commentstyle=\color{codegreen},
    keywordstyle=\color{magenta},
    numberstyle=\tiny\color{codegray},
    stringstyle=\color{codepurple},
    basicstyle=\ttfamily\footnotesize,
    breakatwhitespace=false,
    breaklines=true,
    captionpos=b,
    keepspaces=true,
    numbers=left,
    numbersep=5pt,
    showspaces=false,
    showstringspaces=false,
    showtabs=false,
    tabsize=2
}
\lstset{style=pythonstyle}

\title{Lecture: Self-Attention and the Transformer Revolution\\
From Sequential Processing to Parallel Computation}
\author{Deep Learning Class}
\date{\today}

\begin{document}
\maketitle

\section{BERT: Bidirectional Encoder Representations from Transformers}\label{sec:bert}
\textbf{What it is.} \;BERT is an \emph{encoder-only Transformer} pre-trained on large unlabeled corpora to produce \emph{contextual} token representations. Unlike left-to-right language models, each token attends to both left and right context (bidirectional), and the same pre-trained encoder can be fine-tuned for many downstream tasks.

\medskip
\noindent\textbf{Input representation.}
\begin{itemize}
  \item \emph{Tokens:} subword units (WordPiece/BPE). Sequence is typically packed as \texttt{[CLS]}~Sentence~A~\texttt{[SEP]}~Sentence~B~\texttt{[SEP]}.
  \item \emph{Embeddings (summed):} token embeddings $+$ \emph{learned absolute} position embeddings $+$ segment (token-type) embeddings (A/B) to mark sentence membership.
  \item \emph{[CLS]:} special first token; its final hidden state is used as a sequence summary for classification.
\end{itemize}

\noindent\textbf{Architecture.}
\begin{itemize}
  \item Stack of Transformer encoder blocks: Multi-Head Self-Attention $\rightarrow$ Add\&Norm $\rightarrow$ Positionwise FFN $\rightarrow$ Add\&Norm.
  \item \emph{BERT-Base:} 12 layers, hidden size $d_{\text{model}}=768$, 12 heads ($\sim$110M params). \quad
        \emph{BERT-Large:} 24 layers, $d_{\text{model}}=1024$, 16 heads ($\sim$340M params).
\end{itemize}

\noindent\textbf{Pretraining objectives.}
\begin{enumerate}
  \item \emph{Masked Language Modeling (MLM):} Randomly select $15\%$ of tokens. Of these, replace $80\%$ with \texttt{[MASK]}, $10\%$ with a random token, and keep $10\%$ unchanged. Train to recover the original tokens at masked positions via cross-entropy:
  \[
    \min_{\theta}\; -\sum_{i\in \mathcal{M}} \log p_{\theta}\!\left(x_i \mid x_{\setminus i}\right),
  \]
  where $\mathcal{M}$ indexes masked positions and $x_{\setminus i}$ denotes the full sequence with the $i$-th token hidden.
  \item \emph{Next Sentence Prediction (NSP):} Given a pair \texttt{(A,B)}, predict whether $B$ is the true next sentence after $A$ (half of pairs are replaced with a random $B$). Use the final \texttt{[CLS]} state for a binary classifier.
\end{enumerate}

\noindent\textbf{Why this works (intuitions).}
\begin{itemize}
  \item \emph{Bidirectional context:} MLM forces every position to integrate information from both sides, yielding richer semantics than uni-directional LMs for understanding tasks.
  \item \emph{Robust representations:} The 80/10/10 masking policy discourages over-reliance on the \texttt{[MASK]} token and teaches BERT to detect tokens that ``don’t fit'' their context.
  \item \emph{Discourse signal:} NSP adds a coarse objective for inter-sentence coherence, useful for pairwise tasks (e.g., NLI, QA over a passage given a question).
\end{itemize}

\noindent\textbf{Data scale.} Pretrained on large corpora (e.g., BooksCorpus + Wikipedia, billions of words). Scale exposes the encoder to diverse syntax/semantics and broad world-knowledge cues.

\medskip
\noindent\textbf{Fine-tuning pattern.}
\begin{itemize}
  \item \emph{Sequence classification (e.g., sentiment):} feed \texttt{[CLS]} into a small MLP $\rightarrow$ softmax.
  \item \emph{Token labeling (e.g., NER, POS):} apply a classifier to each token’s final hidden state.
  \item \emph{Extractive QA:} learn start/end classifiers over token positions; predict spans within the passage.
\end{itemize}

\noindent\textbf{Common clarifications.}
\begin{itemize}
  \item Original BERT uses \emph{learned absolute} position embeddings (not sinusoidal).
  \item BERT is not a left-to-right generative LM; it excels at \emph{understanding} tasks. Decoder LMs (e.g., GPT) are preferred for generation.
  \item Subword tokenization (WordPiece/BPE) gives an open vocabulary with manageable size; rare words decompose into frequent subunits.
\end{itemize}

\noindent\textbf{Mental model.} View BERT as a deep, bidirectional text \emph{encoder}. MLM teaches recovery of missing words from full context; NSP adds a simple discourse check. After pretraining, attach a light task head and fine-tune end-to-end.


\end{document}
